\section{Discussion}
\label{sec:discussion}

The results support two findings of different strength. The validity of a test that uses the logged design follows from
its construction and was confirmed on placebo logs, on null outcomes and in controlled environments with known causal
graphs. The finding that a valid map did not improve safe conflict mitigation, in contrast, rests on a single simulated
scenario and a single referee rule.

\subsection{Implications for causal discovery}
\label{sec:disc_discovery}\claim{EMP-POWER}

The validity of the test derives from the logged design rather than from a model of the network, which makes it the
property most likely to transfer to other networks. The randomisation distribution is generated from the logged
propensities alone, whether uniform or state-dependent, so the guarantee holds for any statistic, weighting or prior
direction fixed before each unit draws its mode. The controlled environments confirm that this guarantee belongs to the
design rather than to PMRT, since partial correlation adjusted for the design variables also retained its nominal rate
under a logged state-dependent policy.

Statistical power, in contrast, depends on the choice of statistic. Design-adjusted partial correlation had higher recall in the
latent-confounding environment, whereas PMRT had higher recall on the uniform-randomised simulator logs. PMRT's
matched filter and weights therefore concern power rather than the source of validity.

\subsection{Implications for conflict mitigation}
\label{sec:disc_mitigation}

In the evaluated scenario, mitigation was limited by the guard constraints and objective, not by map accuracy. With
the gate, the PMRT referee made the same decisions as the never-sleep rule, and the oracle map performed no better.
Further improvements in discovery would therefore not have increased the benefit of safe mitigation here. Conversely,
the ungated result shows that one missed edge was sufficient for the referee to breach a guard limit. A map should therefore not drive
mitigation without a step that certifies the effect of each decision on every guard KPI.

\subsection{Deployment requirements}
\label{sec:disc_deploy}\claim{EMP-X3}\claim{EMP-S1-STRICT}\claim{EMP-X4}

The test requires the conflict-mitigation layer to draw each accept-or-reject decision with a known propensity and
record it exactly. The policy may change over time, provided that each recorded propensity is the one used for the
draw. Otherwise the design is misspecified: in the controlled environments, PMRT-GBM and PMRT-Lin exceeded the
nominal null-edge rate when the assumed dither width was 0.8 times the true width, and every PMRT variant did so when
the assumed logging-policy strength was twice its true value (\esm{esm:envs}). The RIC should therefore record the
propensity at the moment of the draw, from the code that performs it, and keep propensities bounded away from 0 and 1.

Randomised logging imposes policy and evidence costs. Random rejection reduced recovery below accept-all because it
rejected requests that a fixed rule would accept. Evidence also accumulates at the rate of decision units, each a 60~s
mode followed by a 90~s outcome window, rather than at the RIC control-loop rate.
The smallest simulator samples, 60 episodes of 600~s, therefore represent hours of network time. At that size, small
p-values occurred at their nominal rate under a stricter significance threshold (\esm{esm:discovery}), although the
small-effect knockout nulls could not exclude a rate up to three times nominal. The logging time required in a deployed
RIC will depend on traffic and decision rates not examined here.

For operators applying learned discovery methods to designed logs, a placebo check provides an inexpensive screening
step. The pipeline is applied unchanged to an additional column drawn from the logged design, which by construction
affects no KPI. Any edge declared from that column disqualifies the pipeline's declarations as causal evidence. Passing
the check is not sufficient; positive causal evidence still requires validity derived from the logging design.

\subsection{Limitations}
\label{sec:disc_limits}\claim{EMP-X15}\claim{EMP-X16}

The simulator results come from one scenario, one incumbent policy and one referee rule, and none involves a real RIC,
a testbed or a compliance claim. The controlled environments model no deployed RAN, and no method evaluated in them
uses the matched filter and weighted BY procedure of simulator PMRT. Their claims apply to the tested
dithered-setpoint design, and their evidence against design-blind tests rests in one environment on a placebo edge
only. Results from one block of seeds also carry seed-level heterogeneity larger than their intervals suggest
(\esm{esm:decision}).

PMRT is asymptotic rather than exact, even on placebo logs. Its FDR guarantee concerns sharp-null declarations, not
signs, relation attribution or knockout materiality labels. The ground truth
targets $\tau_{\mathrm{AA}}$ rather than $\tau_{\pi}$, and its labels depend on $\delta_k$. The design slope and gate
bounds use an approximately predictable outcome centre; replacing it with a strictly predictable centre changed the
map slopes by \val{x16.map.rel_min} to \val{x16.map.rel_max}~\% without changing any sign, moved the bounds by a
median \val{x16.gate.median_pct}~\% of their tolerance and altered no certification (\esm{esm:decision}). No guard
covers the protected slice, and the associational maps use independent-error bounds that may be too narrow for
dependent KPI series by an amount not quantified here.
